\documentclass[11pt,a4paper]{article}

% ---------- Packages ----------
\usepackage[utf8]{inputenc}
\usepackage[T1]{fontenc}
\usepackage[margin=2.2cm]{geometry}
\usepackage{amsmath,amssymb,amsthm}
\usepackage{booktabs}
\usepackage{array}
\usepackage{graphicx}
\usepackage{xcolor}
\usepackage{enumitem}
\usepackage{hyperref}
\usepackage{listings}
\usepackage{algorithm}
\usepackage{algpseudocode}
\usepackage{caption}
\usepackage{titlesec}
\usepackage{fancyhdr}
\usepackage[numbers,sort&compress]{natbib}

% ---------- Colors ----------
\definecolor{navy}{HTML}{1A1A2E}
\definecolor{steel}{HTML}{2D4A7A}
\definecolor{lightgray}{HTML}{F4F4F4}

% ---------- Header / Footer ----------
\pagestyle{fancy}
\fancyhf{}
\fancyhead[L]{\small\textcolor{steel}{Flexible Flow Line Scheduling}}
\fancyhead[R]{\small\textcolor{steel}{Pinedo Ch.~16.3}}
\fancyfoot[C]{\small\thepage}
\renewcommand{\headrulewidth}{0.5pt}
\renewcommand{\footrulewidth}{0pt}

% ---------- Section formatting ----------
\titleformat{\section}
  {\normalfont\Large\bfseries\color{navy}}
  {\thesection}{0.6em}{}
\titleformat{\subsection}
  {\normalfont\large\bfseries\color{steel}}
  {\thesubsection}{0.6em}{}
\titleformat{\subsubsection}
  {\normalfont\bfseries\color{steel}}
  {\thesubsubsection}{0.6em}{}

% ---------- Listings style ----------
\lstdefinestyle{pseudo}{
  backgroundcolor=\color{lightgray},
  basicstyle=\ttfamily\small,
  breaklines=true,
  frame=single,
  rulecolor=\color{steel},
  framesep=6pt,
  xleftmargin=6pt,
  xrightmargin=6pt,
}

\hypersetup{
  colorlinks=true,
  linkcolor=steel,
  citecolor=steel,
  urlcolor=steel,
}

\newtheorem{definition}{Definition}

\title{\vspace{-1.5cm}
  {\color{steel}\Large Operations Research Project Report}\\[0.4cm]
  {\color{navy}\Huge\bfseries Scheduling of a Flexible Flow Line\\ with Limited Buffers and Bypass}\\[0.3cm]
  {\large Chapter 16.3 --- Pinedo, \textit{Scheduling: Theory, Algorithms, and Systems}}
}
\author{}
\date{\today}

\begin{document}

\maketitle
\thispagestyle{empty}

\begin{abstract}
\noindent
This report analyzes the scheduling problem described in Section~16.3 of
Pinedo's \textit{Scheduling: Theory, Algorithms, and Systems}
\citep{pinedo2016scheduling}: the cyclic scheduling of a flexible flow
line with limited buffers and job bypass, originally addressed at IBM via
the Flexible Flow Line Loading (FFLL) algorithm \citep{ibm_ffll}. We
(1)~formalize the problem and its Graham three-field classification,
(2)~derive and critically assess a mixed-integer programming (MIP)
formulation, (3)~describe the three-phase FFLL heuristic and apply it to
two original numerical examples, and (4)~propose an Iterated Greedy with
Local Search (IG-LS) algorithm inspired by the iterated greedy framework of
\citet{ruiz2007simple}, implement it, and benchmark it against FFLL on
ten randomly generated instances. IG-LS reduces the average optimality gap
from 33.8\% (FFLL) to 15.1\%, reaching the lower bound on two of the ten
instances.
\end{abstract}

\vspace{0.5cm}
\tableofcontents
\newpage

% =====================================================================
\section{Problem Statement and Graham Classification}
% =====================================================================

\subsection{Problem Description}

The problem considers a \textbf{flexible flow line (FFL)} -- a production
environment consisting of multiple stages arranged in series, where each
stage contains one or more machines operating in parallel. The specific
context in Section~16.3 of \citet{pinedo2016scheduling} is the assembly
of printed circuit boards (PCBs), an environment for which the Flexible
Flow Line Loading (FFLL) algorithm was originally designed at IBM
\citep{ibm_ffll}.

The environment has the following defining characteristics:

\begin{description}[leftmargin=2.2cm,style=nextline]
  \item[Jobs] Each job represents a batch of relatively small identical
    items (e.g., PCBs). A set of jobs that together represent one full
    day's production mix is called a \emph{Minimum Part Set} (MPS). The
    manufacturing process is repetitive, so the MPS cycles through the
    system continuously.
  \item[Processing] Each job must be processed at every stage on exactly
    one machine. Not all machines at a stage are necessarily identical --
    some may handle only specific job types.
  \item[Bypass] A job may not require processing at every stage. If a job
    skips a stage, the material handling system transports it past that
    stage. This is modeled by setting the processing time of that job at
    the bypassed stage to zero, $p_{ij}=0$.
  \item[Limited buffers] Between stages, buffer space is finite. If a
    buffer is full, the material handling system halts -- or, if
    recirculation is possible, the job bypasses the stage and recirculates.
  \item[Cyclic schedule] Because production is repetitive, the objective
    is a \emph{cyclic schedule}: the best ordering and timing of jobs
    within one MPS cycle, repeated indefinitely.
\end{description}

\subsection{Objectives and Alternative Objective Functions}

The FFLL algorithm pursues two objectives simultaneously
\citep{pinedo2016scheduling}:

\begin{itemize}
  \item \textbf{Primary objective -- maximize throughput.} In cyclic
    scheduling this is equivalent to minimizing the cycle time of one MPS.
    The cycle time cannot be smaller than the workload of the bottleneck
    machine, so the algorithm seeks to approach this lower bound.
  \item \textbf{Secondary objective -- minimize work-in-process (WIP).}
    Since buffer space is limited, reducing WIP lowers the probability of
    buffer overflow and the resulting blocking.
\end{itemize}

Alternative objective functions that could be substituted for $C_{\max}$
include:

\begin{itemize}
  \item Minimizing weighted tardiness $\sum w_j T_j$ -- relevant when jobs
    carry due dates and differing priorities.
  \item Minimizing total/average flow time $\sum C_j$ -- reducing the
    average time a job spends in the system.
  \item Minimizing the number of tardy jobs $\sum U_j$ -- relevant in
    customer-facing settings.
  \item Minimizing maximum lateness $L_{\max}$ -- relevant when due-date
    violations carry hard consequences.
\end{itemize}

\subsection{Graham Notation}

We use the standard three-field notation $\alpha \mid \beta \mid \gamma$ of
\citet{graham1979optimization}, adopting Pinedo's machine-environment
symbol for this specific class of problems. Pinedo denotes the flexible
flow shop by \textbf{FFc}, where $c$ is the number of stages in series,
each stage being a bank of \emph{parallel machines}
\citep{pinedo2016scheduling}. The classification is therefore:

\begin{equation*}
  \boxed{\;FFc \mid bypass,\; b_i < \infty \mid C_{\max}\;}
\end{equation*}

\begin{table}[h!]
\centering
\begin{tabular}{@{}lll@{}}
\toprule
\textbf{Field} & \textbf{Symbol} & \textbf{Meaning} \\
\midrule
$\alpha$ (machine environment) & $FFc$ & Flexible flow shop: $c$ stages in
  series, each stage a \\
  & & bank of parallel machines \citep{pinedo2016scheduling} \\
$\beta$ (constraints) & $bypass$ & Jobs may skip a stage (zero processing
  time there) \\
 & $b_i < \infty$ & Finite buffer capacity between stages \\
$\gamma$ (objective) & $C_{\max}$ & Minimize makespan / cycle time of the
  MPS \\
\bottomrule
\end{tabular}
\caption{Graham three-field classification of the flexible flow line
  scheduling problem.}
\end{table}

\textbf{Remark.} $FFc$ generalizes both the classical flow shop ($Fm$,
single machine per stage) and the parallel-machine environment
($Pm$, a single stage); it reduces to $Fm$ when every stage has exactly one
machine. Adding $bypass$ means routing is no longer uniform across jobs,
and finite buffers ($b_i<\infty$) introduce \emph{blocking}, a phenomenon
absent from most classical scheduling theory. Even the two-machine flow
shop with blocking is known to be NP-hard
\citep{hall1996machine}, and the combination of parallel machines,
bypass, and finite buffers studied here is at least as hard. This
NP-hardness is the central motivation for the heuristic (rather than exact)
solution approach of Section~3.

% =====================================================================
\section{MIP Formulation and Critical Assessment}
% =====================================================================

Pinedo's chapter does not present an explicit mixed-integer program; we
construct one here that mirrors the three decisions made sequentially by
the FFLL algorithm: machine assignment, job sequencing, and timing.

\subsection{Indices, Parameters and Decision Variables}

\paragraph{Indices and parameters.}
\begin{itemize}
  \item $n$: number of jobs in the MPS, indexed $j=1,\dots,n$.
  \item $m$: total number of machines, indexed $i=1,\dots,m$.
  \item $S$: number of stages, indexed $s=1,\dots,S$.
  \item $M_s$: set of machines belonging to stage $s$.
  \item $p_{ij}$: processing time of job $j$ on machine $i$ ($0$ if job $j$
    bypasses the stage containing machine $i$).
  \item $b_s$: buffer capacity between stage $s$ and stage $s+1$.
  \item $L$: a sufficiently large constant (big-$M$).
\end{itemize}

\paragraph{Decision variables.}
\begin{table}[h!]
\centering
\begin{tabular}{@{}lll@{}}
\toprule
Variable & Domain & Meaning \\
\midrule
$x_{ij}$ & $\{0,1\}$ & $1$ if job $j$ is assigned to machine $i$ \\
$y_{ijj'}$ & $\{0,1\}$ & $1$ if job $j$ precedes job $j'$ on machine $i$ \\
$C_{ij}$ & $\mathbb{R}_{\ge 0}$ & Completion time of job $j$ on machine $i$ \\
$T$ & $\mathbb{R}_{\ge 0}$ & Cycle time (makespan of one MPS) \\
\bottomrule
\end{tabular}
\end{table}

\subsection{MIP Model}

\begin{align}
  \text{minimize} \quad & T \label{eq:obj}\\[4pt]
  \text{s.t.}\quad
  &\sum_{i \in M_s} x_{ij} = 1
    && \forall j,\; \forall s \tag{C1}\label{eq:c1}\\
  &C_{ij} \ge C_{ij'} + p_{ij} x_{ij} - L(1 - y_{ijj'})
    && \forall i,\; \forall j \neq j' \tag{C2a}\label{eq:c2a}\\
  &C_{ij'} \ge C_{ij} + p_{ij'} x_{ij'} - L\, y_{ijj'}
    && \forall i,\; \forall j \neq j' \tag{C2b}\label{eq:c2b}\\
  &C_{ij} \ge C_{i'j} + p_{ij} x_{ij}
    && \forall j,\, s,\, i \in M_s,\, i' \in M_{s-1} \tag{C3}\label{eq:c3}\\
  &\Bigl|\{\, j : C_{i'j} \le t < C_{ij} \,\}\Bigr| \le b_s
    && \forall s,\; \forall t \tag{C4}\label{eq:c4}\\
  &T \ge C_{ij}
    && \forall i,j : x_{ij}=1 \tag{C5}\label{eq:c5}\\
  &T \ge W_i = \sum_{j=1}^{n} p_{ij}
    && \forall i \tag{C6}\label{eq:c6}\\
  &x_{ij},\, y_{ijj'} \in \{0,1\},\quad C_{ij}, T \ge 0 \tag{C7}\label{eq:c7}
\end{align}

\noindent
\textbf{(C1)} assigns every job to exactly one machine per stage --
bypassed stages satisfy this trivially since $p_{ij}=0$ for all $i \in
M_s$. \textbf{(C2a--C2b)} are the classical disjunctive constraints
preventing two jobs from overlapping on the same machine. \textbf{(C3)}
enforces the series structure of the flow line: a job cannot start
processing at stage $s$ before completing stage $s-1$. \textbf{(C4)}
restricts the number of jobs simultaneously resident in the buffer between
consecutive stages to the buffer capacity $b_s$; in practice this requires
either time-indexed binary variables or an auxiliary counting formulation,
since it is not linear as written. \textbf{(C5)} defines the cycle time as
the latest completion time across all jobs and machines. \textbf{(C6)} is a
valid inequality stating that the cycle time cannot be smaller than the
total workload of any single machine -- this tightens the LP relaxation
considerably.

\subsection{Model Size and Complexity}

\begin{table}[h!]
\centering
\begin{tabular}{@{}lc@{}}
\toprule
Component & Count \\
\midrule
Binary variables $x_{ij}$ & $m \times n$ \\
Binary sequencing variables $y_{ijj'}$ & $m \times n \times (n-1)$ \\
Continuous variables $C_{ij}$ & $m \times n$ \\
Disjunctive constraints \eqref{eq:c2a}--\eqref{eq:c2b} & $O(m n^2)$ \\
Buffer constraints \eqref{eq:c4} & $O(S \cdot T_{\text{horizon}})$ \\
\bottomrule
\end{tabular}
\caption{Growth of the MIP's size with problem dimensions.}
\end{table}

Even for a modest instance with $m=5$ machines and $n=10$ jobs, the number
of binary sequencing variables already reaches $450$, and the disjunctive
constraints yield a notoriously weak LP relaxation -- a well-documented
issue for job-shop-style MIPs \citep{pinedo2016scheduling}. The buffer
constraint \eqref{eq:c4} is the most problematic component, as a faithful
linearization requires time-indexed binary variables, exploding the model
size for any non-trivial planning horizon.

\subsection{Critical Assessment of the MIP}

\paragraph{Strengths.}
\begin{itemize}
  \item Provides an exact representation; guarantees global optimality
    given sufficient solver time.
  \item The bottleneck lower bound \eqref{eq:c6} tightens the relaxation
    and accelerates branch-and-bound.
  \item Flexible: alternative objectives (WIP, weighted tardiness) can be
    substituted directly into \eqref{eq:obj}.
  \item Useful as an exact benchmark for evaluating heuristic solution
    quality on small instances.
\end{itemize}

\paragraph{Weaknesses.}
\begin{itemize}
  \item The disjunctive constraints \eqref{eq:c2a}--\eqref{eq:c2b} are the
    classical weakness of job-shop MIPs, producing enormous branch-and-bound
    trees even for small instances.
  \item The buffer constraints \eqref{eq:c4} are not linear as stated and
    require an expensive time-indexed reformulation.
  \item The cyclic nature of the schedule -- where the end of one MPS
    cycle feeds directly into the start of the next -- is not captured by
    the static formulation above without additional wrap-around
    constraints.
  \item The bypass feature, while trivial to encode via $p_{ij}=0$,
    interacts subtly with the buffer constraints, since bypassing jobs
    still occupy the material handling system.
  \item The model does not scale: industrial PCB assembly lines may involve
    dozens of stages and hundreds of job types per MPS, far beyond the
    practical reach of exact MIP solvers.
\end{itemize}

\noindent \textbf{Overall verdict.} The MIP is theoretically sound and a
useful small-instance benchmark, but it is not a practical solution method
for industrial-scale instances of this problem -- precisely the motivation
behind IBM's development of the FFLL heuristic \citep{ibm_ffll}.

% =====================================================================
\section{The FFLL Algorithm}
% =====================================================================

\subsection{Algorithm Description}

The Flexible Flow Line Loading (FFLL) algorithm \citep{ibm_ffll,
pinedo2016scheduling} solves the problem in three sequential phases, each
addressing one of the three core decisions: \emph{who} processes \emph{what},
\emph{in what order}, and \emph{when}.

\subsubsection*{Phase 1 --- Machine Allocation (LPT Heuristic)}

Each job is assigned to one machine per stage so that workloads across
parallel machines at that stage are as balanced as possible -- balance
directly determines the bottleneck workload and hence the achievable cycle
time. The Longest Processing Time (LPT) heuristic \citep{graham1969bounds}
is applied independently at every stage with multiple machines:

\begin{enumerate}
  \item Sort all jobs in decreasing order of processing time at that stage.
  \item Assign each job, one at a time, to the currently least-loaded
    machine at that stage.
  \item Repeat until all jobs are assigned.
\end{enumerate}

\subsubsection*{Phase 2 --- Sequencing (Dynamic Balancing Heuristic)}

Let $n$ be the number of jobs and $m$ the number of machines. Define

\begin{align*}
  W_i &= \sum_{j=1}^n p_{ij} && \text{workload of machine } i, \\
  W &= \sum_{i=1}^m W_i && \text{total system workload}, \\
  p_k &= \sum_{i=1}^m p_{ik} && \text{system-wide workload of job } k, \\
  o_{ij} &= p_{ij} - p_j \frac{W_i}{W} && \text{overload contribution of job } j \text{ to machine } i, \\
  O_{ij} &= \sum_{k \in S_j} o_{ik} && \text{cumulative overload on machine } i \text{ after position } j,
\end{align*}

where $S_j$ is the set of jobs loaded up to and including position $j$ in
the sequence. The Dynamic Balancing heuristic greedily minimizes

\begin{equation}
  \sum_{i=1}^{m} \sum_{j=1}^{n} \max(O_{ij}, 0)
  \label{eq:overload}
\end{equation}

by, at each position, selecting from the remaining unscheduled jobs the
one that minimizes the resulting total positive cumulative overload.

\subsubsection*{Phase 3 --- Timing of Releases (Bottleneck Anchoring)}

\begin{enumerate}
  \item Identify the bottleneck machine $i^\ast = \arg\max_i W_i$; the
    cycle time lower bound is $T^\ast = W_{i^\ast}$.
  \item Release all jobs into the system as rapidly as possible and
    simulate the resulting schedule.
  \item Fix the start/completion times on the bottleneck machine -- no
    idle time is permitted there.
  \item For machines \emph{upstream} of the bottleneck, delay processing
    as much as possible without altering job sequences, reducing WIP.
  \item For machines \emph{downstream} of the bottleneck, process jobs as
    early as possible, again without altering sequences.
\end{enumerate}

\subsection{Two Original Application Examples}

The following two examples are constructed independently of the textbook
and differ from Pinedo's Example~16.3.1.

\subsubsection*{Example A --- Two Stages, Four Jobs, No Bypass}

Stage~1 has two parallel machines; Stage~2 has a single machine; the MPS
contains four jobs.

\begin{table}[h!]
\centering
\begin{tabular}{@{}lcccc@{}}
\toprule
 & Job 1 & Job 2 & Job 3 & Job 4 \\
\midrule
Stage 1 & 5 & 8 & 3 & 6 \\
Stage 2 & 4 & 2 & 7 & 5 \\
\bottomrule
\end{tabular}
\caption{Stage processing times $p'_{kj}$ for Example A.}
\end{table}

\textbf{Phase 1 (LPT at Stage 1).} Sorting descending: Job~2~(8), Job~4~(6),
Job~1~(5), Job~3~(3). Assignment: Job~2$\to$M1 (load 8), Job~4$\to$M2 (load
6), Job~1$\to$M2 (load 11), Job~3$\to$M1 (load 11). Both Stage~1 machines
balance perfectly at load 11; Machine~3 (Stage~2) processes all jobs,
$W_3 = 18$. \textbf{Bottleneck: Machine 3, $T^\ast = 18$.}

\textbf{Phase 2 (Dynamic Balancing).} With $W=(11,11,18)$, $\sum W = 40$,
and $p_k=(9,10,10,11)$, evaluating $o_{ij}$ and applying the greedy rule
position-by-position yields the sequence $\mathbf{1 \to 3 \to 2 \to 4}$.

\textbf{Phase 3 (Timing).} The bottleneck machine processes jobs
back-to-back in this order; the achieved cycle time is $T = 18$, matching
the lower bound exactly (optimal for this instance).

\subsubsection*{Example B --- Three Stages, Five Jobs, with Bypass}

Stage~1 has one machine; Stage~2 has two parallel machines; Stage~3 has one
machine; the MPS contains five jobs, with Jobs~2 and~5 bypassing Stage~2.

\begin{table}[h!]
\centering
\begin{tabular}{@{}lccccc@{}}
\toprule
 & Job 1 & Job 2 & Job 3 & Job 4 & Job 5 \\
\midrule
Stage 1 & 3 & 5 & 2 & 4 & 6 \\
Stage 2 & 4 & 0 & 6 & 3 & 0 \\
Stage 3 & 2 & 3 & 1 & 4 & 2 \\
\bottomrule
\end{tabular}
\caption{Stage processing times $p'_{kj}$ for Example B (Jobs 2, 5 bypass Stage 2).}
\end{table}

\textbf{Phase 1 (LPT at Stage 2).} Only Jobs~1, 3, 4 require Stage~2.
Sorting descending: Job~3~(6), Job~1~(4), Job~4~(3). Assignment:
Job~3$\to$M2 (load 6), Job~1$\to$M3 (load 4), Job~4$\to$M3 (load 7). The
resulting workloads are $W=(20,6,7,12)$, $\sum W = 45$.
\textbf{Bottleneck: Machine 1 (Stage 1), $T^\ast = 20$.}

\textbf{Phase 2 (Dynamic Balancing).} With $p_k = (9,8,9,11,8)$, evaluating
$o_{ij}$ and applying the greedy rule position-by-position yields the
sequence $\mathbf{2 \to 1 \to 3 \to 5 \to 4}$.

\textbf{Phase 3 (Timing).} The bottleneck Machine~1 processes jobs
back-to-back in this order; Jobs~2 and~5 bypass Stage~2 directly. The
achieved cycle time is $T = 20$, again matching the lower bound exactly.

\subsection{Critical Assessment of the FFLL Heuristic}

\paragraph{Strengths.}
\begin{itemize}
  \item Computationally efficient: Phase~1 runs in $O(n \log n)$ per
    stage, Phase~2 in $O(n^2 m)$, and Phase~3 in linear time -- tractable
    for the hundreds of job types found in industrial MPS instances.
  \item The three-phase decomposition mirrors the natural decision
    hierarchy: allocation must precede sequencing (workloads must be
    known first), and sequencing must precede timing.
  \item The Dynamic Balancing objective \eqref{eq:overload} has a clear
    theoretical motivation: minimizing burst loading directly reduces
    buffer occupancy and blocking probability.
  \item Bypass is handled naturally, since bypassed stages simply
    contribute zero to all $o_{ij}$ terms.
  \item Extensive practical experiments reported by Pinedo confirm FFLL's
    value as a scheduling tool in real assembly environments
    \citep{pinedo2016scheduling}.
\end{itemize}

\paragraph{Limitations.}
\begin{itemize}
  \item The three phases are solved \emph{independently and sequentially}
    with no feedback loop -- a poor Phase~1 allocation cannot be corrected
    by Phases~2 or~3.
  \item LPT does not guarantee optimal balance; Pinedo himself notes this
    explicitly for Stage~3 of Example~16.3.1.
  \item The Dynamic Balancing heuristic is \emph{greedy with no lookahead}:
    a locally optimal choice at position $j$ may preclude a substantially
    better global sequence.
  \item WIP minimization is addressed only indirectly, through Phase~3's
    timing adjustments, with no explicit guarantee against blocking under
    tight buffers.
  \item The algorithm assumes a fixed, fully known MPS and offers no
    mechanism for incremental re-optimization under dynamic arrivals or
    uncertain processing times.
\end{itemize}

% =====================================================================
\section{Improved Algorithm: Iterated Greedy with Local Search (IG-LS)}
% =====================================================================

\subsection{Design Rationale}

The principal weakness identified in Section~3.3 is that FFLL's Phase~2
is a \emph{non-backtracking greedy} procedure: once a job is placed,
it remains fixed regardless of how later choices unfold. This is precisely
the failure mode that the \emph{iterated greedy} (IG) metaheuristic
framework of \citet{ruiz2007simple} was designed to address for
permutation flow shop scheduling. Since its introduction, IG has been
successfully adapted to numerous structurally similar flow shop variants
-- including blocking flow shops, no-idle flow shops, and sequence-dependent
setup-time problems -- frequently outperforming more elaborate metaheuristics
while remaining simple to implement
\citep{ruiz2007simple,pan2014effective}. This motivates adapting the
same destruction--reconstruction principle to the FFLL sequencing phase.

We propose \textbf{IG-LS}, which retains FFLL's Phase~1 (LPT allocation)
and Phase~3 (bottleneck timing) unchanged, and replaces Phase~2 with three
layers:

\begin{enumerate}
  \item \textbf{Warm start.} Initialize with the FFLL Dynamic Balancing
    sequence.
  \item \textbf{Local search (2-opt).} Exhaustively evaluate pairwise job
    swaps; accept any swap that reduces the \emph{actual simulated
    makespan}; repeat until no improving swap remains.
  \item \textbf{Iterated perturbation.} Remove a random contiguous block of
    jobs and reinsert it at a random position (destruction--reconstruction,
    following \citet{ruiz2007simple}); re-apply local search; accept the
    new solution only if it improves the incumbent makespan. Repeat for a
    fixed iteration budget.
\end{enumerate}

Crucially, IG-LS optimizes the \emph{true simulated makespan} rather than
the overload proxy \eqref{eq:overload} used by Dynamic Balancing -- a
more direct and unbiased optimization signal.

\subsection{Pseudocode}

\begin{algorithm}[h!]
\caption{IG-LS: Iterated Greedy with Local Search}
\begin{algorithmic}[1]
\Procedure{IG-LS}{instance, n\_iter $=150$, block\_size $=2$}
  \State alloc $\gets$ LPT\_Allocate(instance) \Comment{Phase 1, identical to FFLL}
  \State seq $\gets$ DynamicBalancing(alloc) \Comment{warm start}
  \State (seq, best\_ms) $\gets$ LocalSearch2opt(seq, alloc)
  \For{$k = 1$ \textbf{to} n\_iter}
    \State seq$'$ $\gets$ Perturb(seq, block\_size) \Comment{remove \& reinsert block}
    \State (seq$'$, ms$'$) $\gets$ LocalSearch2opt(seq$'$, alloc)
    \If{ms$'$ $<$ best\_ms}
      \State seq $\gets$ seq$'$; \quad best\_ms $\gets$ ms$'$
    \EndIf
  \EndFor
  \State cycle\_time $\gets$ BottleneckTiming(seq, alloc) \Comment{Phase 3, identical to FFLL}
  \State \textbf{return} seq, cycle\_time, best\_ms
\EndProcedure
\end{algorithmic}
\end{algorithm}

\subsection{Experimental Results on Ten Random Instances}

Ten instances were generated with deliberately varied characteristics --
ranging from small 4-job/2-stage problems to larger 12-job/5-stage
problems, with varying bypass probabilities and numbers of machines per
stage (full generator and parameters in the accompanying source code,
\texttt{ffll\_igls.py}). The lower bound (LB) reported is the bottleneck
machine workload after LPT allocation (constraint~\eqref{eq:c6}).

\begin{table}[h!]
\centering
\small
\begin{tabular}{@{}rrrrrrrrr@{}}
\toprule
\# & Jobs & Stg & LB & FFLL & FFLL gap & IG-LS & IG-LS gap & Improv. \\
\midrule
1  & 4  & 2 & 9.0  & 13.00 & 44.4\%  & 12.00 & 33.3\% & 7.7\%  \\
2  & 5  & 2 & 23.0 & 26.00 & 13.0\%  & 25.00 & 8.7\%  & 3.8\%  \\
3  & 6  & 3 & 32.0 & 35.00 & 9.4\%   & \textbf{32.00} & \textbf{0.0\%} & 8.6\%  \\
4  & 7  & 3 & 38.0 & 42.00 & 10.5\%  & \textbf{38.00} & \textbf{0.0\%} & 9.5\%  \\
5  & 8  & 4 & 21.0 & 43.00 & 104.8\% & 34.00 & 61.9\% & 20.9\% \\
6  & 8  & 3 & 34.0 & 49.00 & 44.1\%  & 35.00 & 2.9\%  & 28.6\% \\
7  & 10 & 4 & 46.0 & 56.00 & 21.7\%  & 50.00 & 8.7\%  & 10.7\% \\
8  & 10 & 5 & 49.0 & 56.00 & 14.3\%  & 54.00 & 10.2\% & 3.6\%  \\
9  & 12 & 4 & 66.0 & 97.00 & 47.0\%  & 72.00 & 9.1\%  & 25.8\% \\
10 & 12 & 5 & 63.0 & 81.00 & 28.6\%  & 73.00 & 15.9\% & 9.9\%  \\
\midrule
\multicolumn{5}{r}{\textbf{Average}} & \textbf{33.8\%} & & \textbf{15.1\%} & \textbf{12.9\%} \\
\bottomrule
\end{tabular}
\caption{FFLL vs.\ IG-LS on ten randomly generated instances. Gaps are
  relative to the bottleneck lower bound (LB). Bold entries indicate
  proven-optimal solutions (gap $=0\%$).}
\label{tab:results}
\end{table}

IG-LS improved \emph{every one} of the ten instances. On Instances~3
and~4 it reached the theoretical lower bound, proving optimality. The
largest gains occur on Instances~5, 6, and~9, where FFLL's greedy
sequencing committed to particularly poor early choices that IG-LS's
perturbation and local search were able to correct. Across all ten
instances, the average gap to the lower bound fell from $33.8\%$ (FFLL) to
$15.1\%$ (IG-LS), while all runs completed in well under one second.

\subsection{Critical Assessment of IG-LS}

\paragraph{Strengths over FFLL.}
\begin{itemize}
  \item More than halves the average optimality gap (33.8\% $\to$ 15.1\%).
  \item Reaches proven optimality on two of ten instances.
  \item Improves consistently across all tested instance types, not just
    isolated cases.
  \item Optimizes actual makespan directly, rather than an overload proxy.
  \item Remains computationally cheap (sub-second per instance at this
    scale), consistent with reports in the IG literature that the
    framework is simple yet highly competitive against more elaborate
    metaheuristics \citep{ruiz2007simple}.
\end{itemize}

\paragraph{Remaining limitations.}
\begin{itemize}
  \item The gap to the lower bound remains substantial on some instances
    (e.g., Instance~5: 61.9\%); note the bound itself may be loose, so the
    true optimality gap could be smaller than reported.
  \item The purely greedy acceptance criterion (accept only strict
    improvements) may miss solutions reachable only via temporary uphill
    moves; a simulated-annealing-style acceptance criterion, as used in
    several IG variants \citep{ruiz2007simple}, could help.
  \item The perturbation block size is fixed at~2; an adaptive schedule
    could better balance exploration and exploitation.
  \item Like FFLL, the three phases remain solved independently; a fully
    integrated metaheuristic spanning allocation and sequencing could
    yield further gains at higher computational cost.
  \item Scalability beyond the tested range (up to 12 jobs) is untested;
    the $O(n^2)$ per-iteration cost of 2-opt local search will eventually
    dominate for much larger MPS instances.
\end{itemize}

\bigskip
\noindent\textit{The complete Python implementation (FFLL, IG-LS, instance
generator, and experiment runner) accompanies this report as
\texttt{ffll\_igls.py}.}

\newpage
\bibliographystyle{plainnat}
\bibliography{references}

\end{document}
